\documentclass[12pt,a4paper]{article}

% ── Packages ──────────────────────────────────────────────────────────────────
\usepackage[T1]{fontenc}
\usepackage[utf8]{inputenc}
\usepackage[margin=2.5cm]{geometry}
\usepackage{amsmath,amssymb}
\usepackage{booktabs}
\usepackage{hyperref}
\usepackage{xcolor}
\usepackage{listings}
\usepackage{graphicx}
\usepackage{titlesec}
\usepackage{enumitem}
\usepackage{fancyhdr}
\usepackage{parskip}
\usepackage{microtype}
\usepackage{caption}
\usepackage{longtable}
\usepackage{array}

% ── Hyperref setup ────────────────────────────────────────────────────────────
\hypersetup{
    colorlinks=true,
    linkcolor=blue!70!black,
    urlcolor=blue!70!black,
    citecolor=green!50!black,
    pdftitle={NINA GRB Plugin Technical Documentation},
    pdfauthor={Halla Daas, Insiyah Zujar, Ezzeddin Diab, Qusai Al-Tah}
}

% ── Code listing styles ───────────────────────────────────────────────────────
\definecolor{codebg}{rgb}{0.97,0.97,0.97}
\definecolor{codeframe}{rgb}{0.8,0.8,0.8}
\definecolor{keywordcolor}{rgb}{0.0,0.3,0.7}
\definecolor{commentcolor}{rgb}{0.4,0.6,0.4}
\definecolor{stringcolor}{rgb}{0.7,0.1,0.1}

\lstset{
    backgroundcolor=\color{codebg},
    frame=single,
    rulecolor=\color{codeframe},
    basicstyle=\ttfamily\small,
    keywordstyle=\color{keywordcolor}\bfseries,
    commentstyle=\color{commentcolor}\itshape,
    stringstyle=\color{stringcolor},
    numbers=left,
    numberstyle=\tiny\color{gray},
    numbersep=8pt,
    breaklines=true,
    showstringspaces=false,
    tabsize=4,
    captionpos=b
}

\lstdefinestyle{python}{
    language=Python,
    morekeywords={import,from,as,def,return,class,if,else,elif,for,while,
                  True,False,None,with,try,except,raise,pass,in,not,and,or}
}

\lstdefinestyle{csharp}{
    language={[Sharp]C},
    morekeywords={var,async,await,string,bool,int,double,float,void,new,
                  public,private,protected,static,readonly,override,virtual,
                  using,namespace,class,interface,enum,return,get,set}
}

\lstdefinestyle{bash}{
    language=bash,
    morekeywords={sudo,systemctl,pip,git,cd,mkdir,cp,echo}
}

% ── Header / footer ───────────────────────────────────────────────────────────
\pagestyle{fancy}
\fancyhf{}
\rhead{\small NINA GRB Plugin --- Technical Documentation}
\lhead{\small AUS Senior Design Project, Spring 2026}
\rfoot{\small Page \thepage}
\renewcommand{\headrulewidth}{0.4pt}

% ── Title data ────────────────────────────────────────────────────────────────
\title{
    \vspace{-1cm}
    {\large American University of Sharjah}\\[0.4cm]
    {\large Department of Computer Science and Engineering}\\[1.2cm]
    {\LARGE \textbf{NINA GRB Plugin}}\\[0.3cm]
    {\Large An Autonomous Framework for Transient Alert\\
    Classification and Telescope Image Processing}\\[0.5cm]
    {\large Technical Documentation}\\[0.8cm]
    \rule{\textwidth}{0.5pt}
}
\author{
    Halla Daas\quad Insiyah Zujar\quad Ezzeddin Diab\quad Qusai Al-Tah\\[0.3cm]
    \small Advisors: Dr.\ Tamer Shanableh \& Dr.\ Nidhal Guessoum
}
\date{Spring 2026}

% ══════════════════════════════════════════════════════════════════════════════
\begin{document}
\maketitle
\thispagestyle{empty}
\newpage

\tableofcontents
\newpage

% ══════════════════════════════════════════════════════════════════════════════
\section{Abstract}

This document describes the technical implementation of an autonomous
Gamma-Ray Burst (GRB) follow-up system built as an open-source plugin for
Nighttime Imaging `N' Astronomy (N.I.N.A). The system monitors the
Gamma-ray Coordinates Network (GCN) for new alert emails, uses GPT-4o-mini
to extract astrophysical parameters, filters candidates against observatory
constraints, and autonomously triggers telescope slewing, plate-solving, and
image capture. Captured FITS files are subsequently processed by a Python
photometry pipeline that performs forced aperture photometry, calibrates
against Gaia DR3, constructs a per-event light curve including 3$\sigma$
upper limits, and fits a power-law temporal decay model
$F(t) \propto t^{-\alpha}$. All data flows through Google Firestore, enabling
real-time coordination between the cloud-hosted alert pipeline, the NINA
plugin running on the observatory PC, and the image analysis service.

% ══════════════════════════════════════════════════════════════════════════════
\section{System Overview}

\subsection{End-to-End Pipeline}

\begin{enumerate}[leftmargin=*]
    \item \textbf{Alert Ingestion:} \texttt{gmail\_listener.py} monitors a
          Gmail inbox 24/7 for new GCN Circular emails (designed for
          deployment as a Linux systemd service on a cloud VM).
    \item \textbf{LLM Parsing:} \texttt{parser.py} sends the raw email text
          to the OpenAI GPT-4o-mini API. The model extracts: GRB name, RA,
          Dec, positional error radius, trigger time, space telescope,
          count rate, flux, SNR, and optical magnitude.
    \item \textbf{Firestore Storage:} \texttt{firebase\_client.py} writes
          parsed alerts to the \texttt{grb\_alerts} Firestore collection.
    \item \textbf{Observability Filter:} The NINA plugin's
          \texttt{GRBObservabilityService} reads new alerts and checks
          altitude, declination range, localization uncertainty, GRB age,
          count rate, flux, SNR, moon constraints, and space telescope
          filters. Passing GRBs are written to \texttt{observable\_list}.
    \item \textbf{NINA Trigger:} \texttt{GRBAlertTrigger} polls
          \texttt{observable\_list} every 30 seconds and sets
          \texttt{GRBPendingState} when a new observable GRB is found.
    \item \textbf{Observatory Preparation:} \texttt{GRBObservatoryStateService}
          checks the safety monitor (polls if unsafe or disconnected), shutter
          state (opens if closed), telescope park state (unparks), and camera
          busy state before proceeding.
    \item \textbf{Capture:} \texttt{GRBCaptureInstruction} slews the mount
          (up to 3 retry attempts), runs a plate-solve centering loop
          (up to 3 iterations, 30\arcsec\ threshold), and captures
          $N \times M$\,s LIGHT frames. Custom FITS headers
          (\texttt{GRBNAME}, \texttt{GRB\_RA}, \texttt{GRB\_DEC},
          \texttt{GRB\_ERR}, \texttt{GRBTEL}, \texttt{GRBTRIG}) are injected
          into every saved file.
    \item \textbf{Firestore Notification:} \texttt{FirestoreCapturePoster}
          writes a \texttt{grb\_captures} document with
          \texttt{status = pending\_analysis}, the FITS directory path, and
          the trigger time.
    \item \textbf{Image Analysis:} \texttt{image\_analyzer.py} (running
          continuously on the observatory PC) detects the new document within
          30 seconds and runs the full photometry pipeline (see
          Section~\ref{sec:science}).
    \item \textbf{Results:} The \texttt{grb\_captures} document is updated
          with \texttt{status = analyzed}, the full per-frame light curve,
          decay index $\alpha$, and summary statistics.
\end{enumerate}

\subsection{Architecture Diagram}

\begin{verbatim}
  [Cloud VM - Linux, 24/7 target deployment]
    gmail_listener.py  -->  parser.py (GPT-4o-mini)
                                  |
                             Firestore
                          grb_alerts collection
                                  |
  [Observatory Windows PC]        |
    GRBAlertTrigger  <-----------/
    (polls every 30s)
           |
    GRBObservabilityService
    GRBObservatoryStateService
           |
    GRBCaptureInstruction
    (slew -> plate-solve -> capture)
           |
    Firestore: grb_captures
    status = pending_analysis
           |
    image_analyzer.py  (24/7 on observatory PC)
           |
    Firestore: grb_captures
    status = analyzed + light_curve + decay_alpha
\end{verbatim}

% ══════════════════════════════════════════════════════════════════════════════
\section{Installation \& Requirements}

\subsection{Python Dependencies}

\begin{lstlisting}[style=bash, caption={Install Python dependencies}]
pip install firebase-admin openai astropy astroquery photutils \
            scipy numpy rawpy astral defusedxml
\end{lstlisting}

\subsection{C\# / NINA Plugin Dependencies}

\begin{itemize}
    \item .NET 8 SDK
    \item N.I.N.A.\ Plugin SDK (NuGet)
    \item \texttt{Google.Apis.Auth} (NuGet)
    \item \texttt{Newtonsoft.Json} (NuGet)
\end{itemize}

\subsection{Firebase / Firestore Setup}

\begin{enumerate}
    \item Create a Firebase project at \url{https://console.firebase.google.com}.
    \item Enable Firestore in Native mode.
    \item Generate a service account key:
          \textit{Project Settings $\to$ Service Accounts $\to$
          Generate new private key}.
    \item Save the JSON file as
          \texttt{GCN\_parser/firebase\_service\_account.json}.
    \item The C\# plugin reads the same file from the plugin output
          directory (auto-copied via PostBuild xcopy).
\end{enumerate}

\subsection{OpenAI API Key}

Create \texttt{GCN\_parser/.env} with:

\begin{lstlisting}[style=bash]
OPENAI_API_KEY=sk-proj-...your-key-here...
\end{lstlisting}

\subsection{Gmail OAuth Setup}

\begin{enumerate}
    \item Enable the Gmail API in Google Cloud Console.
    \item Create OAuth 2.0 credentials and download
          \texttt{credentials.json} to \texttt{GCN\_parser/}.
    \item Run \texttt{gmail\_listener.py} once manually to complete the
          OAuth flow; \texttt{token.json} is saved for future runs.
\end{enumerate}

% ══════════════════════════════════════════════════════════════════════════════
\section{Component Reference}

\subsection{\texttt{gmail\_listener.py}}

Monitors a Gmail inbox for new GCN Circular emails using the Gmail API.

\begin{itemize}
    \item Polls the inbox at a configurable interval.
    \item Identifies GCN Circulars by subject line pattern matching.
    \item Passes raw email body text to \texttt{parser.py}.
    \item Marks processed emails as read to avoid reprocessing.
\end{itemize}

\subsection{\texttt{parser.py}}

Extracts structured GRB parameters from raw GCN Circular email text using
the OpenAI GPT-4o-mini API.

\begin{lstlisting}[style=python, caption={Core LLM extraction call}]
response = client.chat.completions.create(
    model="gpt-4o-mini",
    messages=[
        {"role": "system", "content": SYSTEM_PROMPT},
        {"role": "user",   "content": email_body}
    ],
    response_format={"type": "json_object"}
)
\end{lstlisting}

\textbf{Extracted fields:} \texttt{grb\_name}, \texttt{ra\_deg},
\texttt{dec\_deg}, \texttt{error\_arcmin}, \texttt{trigger\_time},
\texttt{telescope}, \texttt{count\_rate}, \texttt{flux}, \texttt{snr},
\texttt{magnitude}.

\subsection{\texttt{firebase\_client.py}}

Writes parsed GRB alert documents to the Firestore \texttt{grb\_alerts}
collection. Uses \texttt{python-dotenv} with script-relative path resolution
to load credentials regardless of the working directory.

\subsection{\texttt{image\_analyzer.py}}

The core photometry service. Polls Firestore every 30 seconds for
\texttt{grb\_captures} documents with \texttt{status = pending\_analysis}
and runs the full analysis pipeline (see Section~\ref{sec:science}).

\textbf{Key constants:}

\begin{center}
\begin{tabular}{lll}
\toprule
Constant & Value & Description \\
\midrule
\texttt{APERTURE\_RADIUS\_PX} & 8.0 px & Photometry aperture radius \\
\texttt{ANNULUS\_INNER\_PX}   & 10.0 px & Background annulus inner radius \\
\texttt{ANNULUS\_OUTER\_PX}   & 15.0 px & Background annulus outer radius \\
\texttt{DETECTION\_SIGMA}     & 3.0 & SNR detection threshold \\
\texttt{SEARCH\_WINDOW\_MIN}  & 15 min & FITS file search window ($\pm$) \\
\texttt{MIN\_FRAMES\_FOR\_FIT}& 3 & Minimum detections for decay fit \\
\texttt{POLL\_INTERVAL\_SECONDS} & 30 s & Firestore polling interval \\
\bottomrule
\end{tabular}
\end{center}

\textbf{Key functions:}

\begin{itemize}
    \item \texttt{find\_fits\_files(fits\_directory, capture\_time)} ---
          Searches recursively for \texttt{*.fits}, \texttt{*.fit},
          \texttt{*.xisf}, \texttt{*.cr2}, \texttt{*.cr3} files within
          $\pm15$ minutes of the recorded capture time.
    \item \texttt{load\_image(path)} --- Loads FITS or Canon RAW files,
          returning \texttt{(data, header)}.
    \item \texttt{get\_obs\_time(fits\_path, header)} --- Reads
          \texttt{DATE-OBS} from the FITS header; falls back to file
          modification time.
    \item \texttt{measure\_magnitude(fits\_path, ra, dec, error\_arcmin,
          trigger\_time)} --- Full per-frame photometry (see
          Section~\ref{sec:science}).
    \item \texttt{fit\_power\_law\_decay(light\_curve)} --- Fits
          $m(t) = m_0 + 2.5\alpha\log_{10}(t)$ via
          \texttt{scipy.optimize.curve\_fit}.
    \item \texttt{compute\_upper\_limit(bkg\_rms, aperture\_r\_px,
          zero\_point)} --- Computes the 3$\sigma$ limiting magnitude.
    \item \texttt{watch\_captures()} --- Main polling loop.
\end{itemize}

\subsection{\texttt{nina\_launcher.py}}

Automatically launches NINA with the nightly GRB sequence at the correct
time each evening.

\begin{lstlisting}[style=python, caption={Sunset calculation with astral}]
from astral import LocationInfo
from astral.sun import sun

location = LocationInfo(
    name="Observatory",
    timezone=OBSERVATORY_TZ,
    latitude=OBSERVATORY_LAT,
    longitude=OBSERVATORY_LON,
)
s = sun(location.observer, date=today, tzinfo=location.timezone)
launch_time = s["sunset"].astimezone(timezone.utc) \
              - timedelta(minutes=MINUTES_BEFORE_SUNSET)
\end{lstlisting}

\textbf{Configuration variables} (edit at top of file):

\begin{itemize}
    \item \texttt{OBSERVATORY\_LAT} / \texttt{OBSERVATORY\_LON} /
          \texttt{OBSERVATORY\_TZ} --- Observatory location
    \item \texttt{MINUTES\_BEFORE\_SUNSET} --- Default: 5 minutes
    \item \texttt{NINA\_EXE} --- Path to \texttt{NINA.exe}
    \item \texttt{SEQUENCE\_FILE} --- Path to the saved
          \texttt{.json} sequence file
\end{itemize}

NINA is launched with:
\begin{lstlisting}[style=bash]
NINA.exe /sequenceFile "C:\GRB\grb_nightly.json" /runSequence
\end{lstlisting}

\subsection{\texttt{setup\_task\_scheduler.ps1}}

PowerShell script (run once as Administrator) that registers three Windows
Task Scheduler tasks:

\begin{center}
\begin{tabular}{lll}
\toprule
Task Name & Trigger & Purpose \\
\midrule
\texttt{grb-nina-launcher} & At startup & Sleeps until sunset$-$5\,min, launches NINA \\
\texttt{grb-nina-launcher-noon} & Daily 12:00 & Safety re-trigger if PC rebooted after sunrise \\
\texttt{grb-image-analyzer} & At startup & Runs \texttt{image\_analyzer.py} 24/7 \\
\bottomrule
\end{tabular}
\end{center}

\begin{lstlisting}[style=bash, caption={One-time setup}]
powershell -ExecutionPolicy Bypass -File setup_task_scheduler.ps1
\end{lstlisting}

\subsection{\texttt{GRBAlertTrigger.cs}}

NINA sequence trigger (\texttt{ISequenceTrigger}) that polls the
\texttt{observable\_list} Firestore collection every 30 seconds. When a
new GRB is found, it stores the event in \texttt{GRBPendingState} so that
the next \texttt{GRBCaptureInstruction} execution picks it up.

\subsection{\texttt{GRBCaptureInstruction.cs}}

NINA sequence item (\texttt{ISequenceItem}) that performs the full
slew-and-capture workflow:

\begin{enumerate}
    \item Reads \texttt{GRBPendingState} and returns immediately if empty.
    \item Calls \texttt{GRBObservatoryStateService.EvaluateAndPrepareAsync()}.
    \item Saves current telescope coordinates for post-capture return slew.
    \item Calls \texttt{SlewAndCenter()} with up to 3 slew retries and
          3 plate-solve centering iterations (30\arcsec\ threshold).
    \item Captures $N \times M$\,s LIGHT frames (configurable via Options).
    \item Calls \texttt{FirestoreCapturePoster.PostCaptureAsync()}.
    \item Marks the GRB as observed via \texttt{GRBObservabilityService
          .MarkAsObserved()}.
    \item Returns the mount to its original position.
\end{enumerate}

\subsection{\texttt{GRBObservatoryStateService.cs}}

Pre-flight safety checks before any GRB capture attempt. Decision tree:

\begin{enumerate}
    \item \textbf{Safety monitor} --- If disconnected or UNSAFE: poll every
          10\,s until connected \textit{and} safe.
    \item \textbf{Shutter closing} --- Wait for fully closed, then open.
    \item \textbf{Shutter closed} --- Open and wait until fully open.
    \item \textbf{Shutter opening} --- Wait until fully open.
    \item \textbf{Safety re-check} after shutter operations --- If UNSAFE:
          close shutter, wait for safe, re-open.
    \item \textbf{Telescope parked} --- Unpark automatically.
    \item \textbf{Camera exposing} --- Wait up to \texttt{MaxWaitMinutes};
          interrupt if threshold exceeded.
\end{enumerate}

\subsection{\texttt{FirestoreCapturePoster.cs}}

Posts capture records to Firestore using the REST API authenticated via a
Google service account (\texttt{GoogleCredential}). Key methods:

\begin{itemize}
    \item \texttt{PostCaptureAsync(grb, exposureCount, exposureSeconds,
          fitsDirectory)} --- Creates a \texttt{grb\_captures} document
          with \texttt{status = pending\_analysis}.
    \item \texttt{PostObservableAlertAsync(...)} --- Upserts the GRB into
          \texttt{observable\_list} (PATCH if exists, POST if new).
    \item \texttt{PostSkippedAsync(grb, reason)} --- Records unobservable
          GRBs with \texttt{status = skipped} and the reason.
\end{itemize}

\subsection{\texttt{GRBObservabilityService.cs}}

Filters incoming GRB alerts against user-configured constraints:

\begin{center}
\begin{tabular}{ll}
\toprule
Filter & Description \\
\midrule
Altitude & Target must be above minimum elevation \\
Declination & Must be within configured Dec range \\
Localization uncertainty & \texttt{error\_arcmin} $\leq$ \texttt{MaxLocalUncertainty} \\
GRB age & Hours since trigger $\leq$ \texttt{GRB\_age} \\
Count rate & $\geq$ configured threshold \\
Flux & $\geq$ configured threshold \\
SNR & $\geq$ configured threshold \\
Moon altitude & Moon must be below \texttt{MaxMoonAlt} \\
Moon separation & Angular distance $\geq$ \texttt{MinMoonSep} \\
Moon phase & Illumination $\leq$ \texttt{MaxMoonPhase} \\
Space telescope & Must match selected instruments \\
\bottomrule
\end{tabular}
\end{center}

% ══════════════════════════════════════════════════════════════════════════════
\section{Scientific Pipeline}
\label{sec:science}

\subsection{Background Subtraction}

Sigma-clipped statistics are computed over the full image to estimate the
sky background level and noise, rejecting bright stars and cosmic rays:

\begin{equation}
    (\mu_\mathrm{sky},\, \sigma_\mathrm{bkg}) =
    \mathrm{SigmaClip}(\mathrm{image},\, \sigma = 3)
\end{equation}

The background-subtracted image is:
\begin{equation}
    I_\mathrm{sub}(x,y) = I(x,y) - \mu_\mathrm{sky}
\end{equation}

\subsection{Forced Aperture Photometry}

Rather than searching for a source, the flux is measured at the exact
reported GRB coordinates regardless of detection status (forced photometry,
following Kumar et al.\ 2025):

\begin{equation}
    F_\mathrm{GRB} = \sum_{(x,y)\,\in\,\mathcal{A}} I_\mathrm{sub}(x,y)
\end{equation}

where $\mathcal{A}$ is a circular aperture of radius $r = 8$\,px centred
on the WCS-projected GRB position.

\subsection{Signal-to-Noise Ratio}

\begin{equation}
    \mathrm{SNR} = \frac{F_\mathrm{GRB}}{\sigma_\mathrm{bkg}
    \sqrt{\pi r^2}}
\end{equation}

\subsection{Gaia DR3 Zero-Point Calibration}

Reference stars within 5\arcmin\ of the GRB position are queried from
Gaia DR3 (Cheng et al.\ 2025). For each reference star $i$:

\begin{equation}
    \mathrm{ZP}_i = m_{\mathrm{Gaia},i} + 2.5 \log_{10}(F_{\mathrm{ref},i})
\end{equation}

The field zero-point is the median over all $N_\mathrm{ref} \geq 3$ stars:

\begin{equation}
    \mathrm{ZP} = \mathrm{median}(\mathrm{ZP}_1, \ldots, \mathrm{ZP}_{N_\mathrm{ref}})
    \qquad \sigma_\mathrm{ZP} = \mathrm{std}(\mathrm{ZP}_1, \ldots,
    \mathrm{ZP}_{N_\mathrm{ref}})
\end{equation}

\subsection{Calibrated Magnitude (Detection, SNR $\geq 3\sigma$)}

\begin{equation}
    m = \mathrm{ZP} - 2.5 \log_{10}(F_\mathrm{GRB})
\end{equation}

Photometric uncertainty (Cheng et al.\ 2025):

\begin{equation}
    \sigma_m = \sqrt{\sigma_\mathrm{Poisson}^2 + \sigma_\mathrm{ZP}^2}
    \qquad \text{where} \quad
    \sigma_\mathrm{Poisson} = \frac{2.5}{\ln 10} \cdot \frac{1}{\mathrm{SNR}}
\end{equation}

\subsection{3$\sigma$ Upper Limit (Non-detection, SNR $< 3\sigma$)}

When no source is confirmed, the 3$\sigma$ limiting magnitude is reported
(Kumar et al.\ 2025):

\begin{equation}
    m_\mathrm{lim} = \mathrm{ZP} - 2.5 \log_{10}
    \!\left(3\,\sigma_\mathrm{bkg}\sqrt{\pi r^2}\right)
\end{equation}

\subsection{Power-Law Temporal Decay}

GRB optical afterglows fade as a power law in time (Del Vecchio et al.\
2016):

\begin{equation}
    F(t) \propto t^{-\alpha}
\end{equation}

In magnitude space this becomes a linear model:

\begin{equation}
    m(t) = m_0 + 2.5\,\alpha\,\log_{10}(t)
\end{equation}

This is fit across all frames with confirmed detections
($\mathrm{SNR} \geq 3$, $\Delta t > 0$) using
\texttt{scipy.optimize.curve\_fit} with photometric uncertainties as
weights. A minimum of 3 detection frames spanning $>3$\,min is required.

Typical values for optical afterglows: $\alpha \approx 0.8$--$1.5$
(Del Vecchio et al.\ 2016, survey of 176 Swift GRBs).

% ══════════════════════════════════════════════════════════════════════════════
\section{Configuration}

All NINA plugin settings are accessible via
\textit{Options $\to$ Plugins $\to$ GRB Helper} in NINA.

\begin{longtable}{lllp{5.5cm}}
\toprule
Setting & Type & Default & Description \\
\midrule
\endfirsthead
\toprule
Setting & Type & Default & Description \\
\midrule
\endhead
\bottomrule
\endfoot
\texttt{DecMin} & double & $-40°$ & Minimum declination \\
\texttt{DecMax} & double & $+90°$ & Maximum declination \\
\texttt{MaxLocalUncertainty} & double & 30 arcmin & Max positional error radius \\
\texttt{GRB\_age} & double & 0 h & Max hours since trigger (0 = disabled) \\
\texttt{CountRate} & double & 0 & Min count rate (counts/s; 0 = disabled) \\
\texttt{Flux\_1}, \texttt{Flux\_2} & float, int & 0 & Min flux mantissa $\times 10^{\text{exp}}$ \\
\texttt{SNR} & double & 0 & Min SNR from alert (0 = disabled) \\
\texttt{Altitude} & double & 30° & Min target altitude above horizon \\
\texttt{Mag} & double & 20 & Max expected optical magnitude \\
\texttt{isSwift} & bool & false & Require Swift detection \\
\texttt{isEP} & bool & false & Require Einstein Probe detection \\
\texttt{isSVOM} & bool & false & Require SVOM detection \\
\texttt{isFermi} & bool & false & Require Fermi detection \\
\texttt{isOther} & bool & false & Allow other telescopes \\
\texttt{ExposureCount} & int & 3 & Number of LIGHT frames per GRB \\
\texttt{ExposureTime} & double & 30 s & Duration of each LIGHT frame \\
\texttt{Binning} & int & 1 & Camera binning ($1=1{\times}1$) \\
\texttt{AstroTwilight} & double & $-12°$ & Max sun altitude for observations \\
\texttt{MaxMoonAlt} & double & 20° & Max moon altitude \\
\texttt{MinMoonSep} & double & 30° & Min moon--GRB angular separation \\
\texttt{MaxMoonPhase} & double & 75\% & Max lunar illumination \\
\texttt{MaxWaitMinutes} & int & 5 min & Max wait if camera is exposing \\
\texttt{ObservatoryCode} & string & --- & MPC observatory code \\
\texttt{SelectedProfile} & string & --- & Active NINA profile name \\
\end{longtable}

% ══════════════════════════════════════════════════════════════════════════════
\section{Deployment Guide}

\subsection{Cloud VM --- Linux (Alert Pipeline)}

The alert ingestion pipeline is designed for deployment as systemd services
on a Linux VM (e.g.\ GCP \texttt{e2-micro} free tier, AWS \texttt{t3.micro},
or DigitalOcean \$4/mo droplet). \textit{Note: currently validated locally;
cloud deployment is the planned production target.}

\begin{lstlisting}[style=bash, caption={Install and clone}]
git clone https://github.com/AUS-Senior-Design/NINA_GRB_Plugin_OpenSource
cd NINA_GRB_Plugin_OpenSource/GCN_parser
pip install -r requirements.txt
\end{lstlisting}

\textbf{Create \texttt{/etc/systemd/system/grb-parser.service}:}

\begin{lstlisting}[style=bash]
[Unit]
Description=GRB GCN Gmail Listener + Parser
After=network.target

[Service]
User=ubuntu
WorkingDirectory=/home/ubuntu/NINA_GRB_Plugin_OpenSource/GCN_parser
ExecStart=/usr/bin/python3 gmail_listener.py
Restart=always
RestartSec=10

[Install]
WantedBy=multi-user.target
\end{lstlisting}

\begin{lstlisting}[style=bash, caption={Enable services}]
sudo systemctl enable grb-parser
sudo systemctl start grb-parser
sudo systemctl status grb-parser
\end{lstlisting}

\subsection{Observatory PC --- Windows}

\subsubsection{Build and Deploy the NINA Plugin}

\begin{enumerate}
    \item Open \texttt{Demo\_2/Demo\_2.sln} in Visual Studio 2022.
    \item Build in \texttt{Debug} or \texttt{Release} mode.
    \item The PostBuild step automatically copies the plugin DLL and
          \texttt{firebase\_service\_account.json} to:\\
          \texttt{\%localappdata\%\textbackslash NINA\textbackslash
          Plugins\textbackslash 3.0.0\textbackslash Sd.NINA.Demo2\textbackslash}
\end{enumerate}

\subsubsection{Configure \texttt{nina\_launcher.py}}

Edit the configuration constants at the top of the file:

\begin{lstlisting}[style=python]
OBSERVATORY_LAT       = 30.0444      # decimal degrees North
OBSERVATORY_LON       = 31.2357      # decimal degrees East
OBSERVATORY_TZ        = "Africa/Cairo"
MINUTES_BEFORE_SUNSET = 5
NINA_EXE  = r"C:\Program Files\N.I.N.A. ...\NINA.exe"
SEQUENCE_FILE = r"C:\GRB\grb_nightly.json"
\end{lstlisting}

\subsubsection{Register Task Scheduler Tasks}

Run once as Administrator:

\begin{lstlisting}[style=bash]
powershell -ExecutionPolicy Bypass -File setup_task_scheduler.ps1
\end{lstlisting}

\subsubsection{NINA Advanced Sequencer Setup}

Build the following sequence in NINA and save to \texttt{SEQUENCE\_FILE}:

\begin{verbatim}
START AREA
  Connect Equipment
  Cool Camera
  Unpark Scope
  Wait If Sun Altitude > -18 deg

TARGET AREA
  [Loop] Wait If Sun Altitude < -18 deg
    TRIGGER: GRBAlertTrigger
    GRBCaptureInstruction
    Wait for Timespan: 00:00:30

END AREA
  Warm Camera
  Park Scope
  Disconnect Equipment
\end{verbatim}

% ══════════════════════════════════════════════════════════════════════════════
\section{Firestore Schema}

\subsection{\texttt{grb\_alerts} Collection}

\begin{center}
\begin{tabular}{lll}
\toprule
Field & Type & Description \\
\midrule
\texttt{grb\_name} & string & GRB designation (e.g.\ \texttt{GRB 260413A}) \\
\texttt{ra\_deg} & number & Right Ascension in degrees (J2000) \\
\texttt{dec\_deg} & number & Declination in degrees (J2000) \\
\texttt{error\_arcmin} & number & Positional uncertainty in arcminutes \\
\texttt{trigger\_time} & string & ISO 8601 UTC trigger timestamp \\
\texttt{telescope} & string & Detecting instrument (Swift, EP, etc.) \\
\texttt{count\_rate} & number & X-ray count rate (counts/s) \\
\texttt{flux} & number & Energy flux (erg\,cm$^{-2}$\,s$^{-1}$) \\
\texttt{snr} & number & Signal-to-noise ratio from alert \\
\texttt{magnitude} & number & Reported optical magnitude \\
\texttt{parsed\_at} & timestamp & Server timestamp of ingestion \\
\bottomrule
\end{tabular}
\end{center}

\subsection{\texttt{observable\_list} Collection}

Same fields as \texttt{grb\_alerts} plus:

\begin{center}
\begin{tabular}{lll}
\toprule
Field & Type & Description \\
\midrule
\texttt{status} & string & \texttt{pending} / \texttt{observed} / \texttt{skipped} \\
\texttt{window\_start} & string & ISO 8601 observability window start \\
\texttt{window\_end} & string & ISO 8601 observability window end \\
\texttt{site\_lat} & number & Observatory latitude \\
\texttt{site\_long} & number & Observatory longitude \\
\texttt{mpc\_code} & string & MPC observatory code \\
\texttt{recorded\_at} & string & ISO 8601 insertion timestamp \\
\bottomrule
\end{tabular}
\end{center}

\subsection{\texttt{grb\_captures} Collection}

\begin{center}
\begin{tabular}{lll}
\toprule
Field & Type & Description \\
\midrule
\texttt{grb\_name} & string & GRB designation \\
\texttt{ra\_deg} & number & Right Ascension (degrees) \\
\texttt{dec\_deg} & number & Declination (degrees) \\
\texttt{error\_arcmin} & number & Positional uncertainty \\
\texttt{trigger\_time} & string & ISO 8601 GRB trigger time \\
\texttt{capture\_time} & string & ISO 8601 NINA capture completion time \\
\texttt{exposure\_count} & number & Number of frames captured \\
\texttt{exposure\_seconds} & number & Exposure duration per frame \\
\texttt{fits\_directory} & string & Local path to FITS files \\
\texttt{status} & string & \texttt{pending\_analysis} / \texttt{analyzed} / \texttt{skipped} \\
\texttt{light\_curve} & array & Per-frame photometry results (see below) \\
\texttt{decay\_alpha} & number & Power-law decay index $\alpha$ \\
\texttt{decay\_alpha\_error} & number & $1\sigma$ uncertainty on $\alpha$ \\
\texttt{decay\_fit\_note} & string & Fit quality description \\
\texttt{magnitude\_mean} & number & Mean calibrated magnitude \\
\texttt{deepest\_upper\_limit} & number & Deepest 3$\sigma$ upper limit \\
\texttt{dss\_peak\_sigma} & number & DSS source significance at GRB position \\
\texttt{known\_sources\_in\_field} & number & SIMBAD object count in error radius \\
\texttt{analyzed\_at} & timestamp & Server timestamp of analysis completion \\
\bottomrule
\end{tabular}
\end{center}

\textbf{Light curve entry fields:}
\texttt{filename}, \texttt{obs\_time}, \texttt{delta\_t\_hours},
\texttt{magnitude}, \texttt{mag\_error}, \texttt{snr}, \texttt{upper\_limit},
\texttt{is\_detection}, \texttt{zero\_point}, \texttt{zp\_scatter},
\texttt{n\_ref\_stars}, \texttt{note}.

% ══════════════════════════════════════════════════════════════════════════════
\section{Running Tests}

\subsection{\texttt{test\_analyzer.py} --- End-to-End Test Harness}

Tests the full image analysis pipeline without a telescope or NINA by
simulating a GRB observation session using real DSS FITS images of the
Crab Nebula:

\begin{lstlisting}[style=python, caption={Run the test harness}]
# Terminal 1 - push a simulated capture to Firestore
cd GCN_parser
python test_analyzer.py

# Terminal 2 - run the analyzer (picks up within 30s)
python image_analyzer.py
\end{lstlisting}

The test harness:
\begin{itemize}
    \item Downloads 4 DSS FITS frames of the Crab Nebula via
          \texttt{astroquery.skyview}.
    \item Injects \texttt{DATE-OBS} headers at simulated offsets:
          $+6$, $+9$, $+12$, $+18$ minutes after a synthetic trigger.
    \item Pushes a \texttt{grb\_captures} document to Firestore with
          \texttt{status = pending\_analysis}.
    \item \texttt{image\_analyzer.py} processes it within 30 seconds and
          writes back the full light curve and decay index.
\end{itemize}

Expected output from \texttt{image\_analyzer.py}:

\begin{lstlisting}[style=bash]
[Analyzer] Found 4 FITS file(s) for GRB_TEST_260413
[Analyzer]   frame_t06min.fits -> DETECTION mag=~16.xx +/- 0.xxx SNR=~xx
[Analyzer]   frame_t09min.fits -> ...
[Analyzer]   frame_t12min.fits -> ...
[Analyzer]   frame_t18min.fits -> ...
[Analyzer] Power-law decay: alpha = x.xxx +/- x.xxx
[Analyzer] Done: GRB_TEST_260413 | detections=4 | upper_limits=0 | alpha=x.xxx
\end{lstlisting}

% ══════════════════════════════════════════════════════════════════════════════
\section{References}

\begin{enumerate}
    \item Kumar, H.\ et al.\ (2025). ``Optical follow-up of gamma-ray burst
          afterglows with the Devasthal Optical Telescope.''
          \textit{Monthly Notices of the Royal Astronomical Society},
          544, 1541.
          DOI: \href{https://doi.org/10.1093/mnras/staf123}{10.1093/mnras/staf123}

    \item Cheng, S.\ et al.\ (2025). ``Optical and near-infrared observations
          of GRB afterglows: Gaia-calibrated photometry and light-curve
          analysis.'' \textit{The Astrophysical Journal}, 979, 38.
          DOI: \href{https://doi.org/10.3847/1538-4357/ad9f12}{10.3847/1538-4357/ad9f12}

    \item Del Vecchio, R.\ et al.\ (2016). ``Probing gamma-ray burst
          afterglow physics with optical flux-density and spectral-index
          light curves.'' \textit{The Astrophysical Journal}, 828, 36.
          DOI: \href{https://doi.org/10.3847/0004-637X/828/1/36}{10.3847/0004-637X/828/1/36}
\end{enumerate}

% ══════════════════════════════════════════════════════════════════════════════
\appendix
\section{Repository Structure}

\begin{verbatim}
NINA_GRB_Plugin_OpenSource/
├── GCN_parser/
│   ├── gmail_listener.py        Gmail watcher for GCN Circulars
│   ├── parser.py                GPT-4o-mini email parser
│   ├── firebase_client.py       Firestore write client
│   ├── image_analyzer.py        Full photometry pipeline (24/7)
│   ├── test_analyzer.py         End-to-end test harness (DSS/Crab)
│   ├── nina_launcher.py         Sunset-aware NINA auto-launcher
│   ├── setup_task_scheduler.ps1 Windows Task Scheduler setup
│   ├── firebase_service_account.json  Firebase credentials
│   └── .env                     OPENAI_API_KEY
│
├── Demo_2/
│   ├── Demo2SequenceItems/
│   │   ├── GRBAlertTrigger.cs            Firestore poller trigger
│   │   ├── GRBCaptureInstruction.cs      Full slew + capture
│   │   └── GRBCaptureOnlyInstruction.cs  Camera-only test mode
│   ├── Services/
│   │   ├── FirestoreCapturePoster.cs     Firestore REST poster
│   │   ├── FirestoreGrbListener.cs       Background listener
│   │   ├── GRBObservatoryStateService.cs Pre-flight checks
│   │   └── GRBObservabilityService.cs    Alert filter logic
│   ├── Models/GRBEvent.cs                GRB data model
│   ├── Options.xaml                      Plugin settings UI
│   └── Demo_2.csproj
│
├── DOCUMENTATION.tex            This document
└── GRB_alerts/                  Archived GRB alert CSV files
\end{verbatim}

\end{document}
